\section{RENCANA DESAIN UI/UX}

\subsection{Pendekatan User Research yang Akan Dilakukan}

Untuk memahami kebutuhan pengguna sebelum mulai implementasi, akan dilakukan penelitian sebagai berikut:

\begin{enumerate}[leftmargin=*]
    \item \textbf{Wawancara dengan UMKM Kuliner} — Akan dilakukan wawancara mendalam untuk memahami tantangan mereka dalam pengelolaan makanan surplus.
    
    \item \textbf{Survei Online} — Akan disebarkan kuesioner kepada konsumen potensial untuk mengetahui preferensi pembelian makanan surplus.
    
    \item \textbf{Observasi Lapangan} — Akan dilakukan pengamatan langsung terhadap operasional UMKM kuliner untuk mengidentifikasi pain points.
\end{enumerate}

\textbf{CATATAN}: Research ini akan dilakukan pada fase awal development, belum dilaksanakan saat ini.

\subsection{Target User Persona (Draft)}

\subsubsection{Persona 1: UMKM Kuliner (Draft)}
\begin{table}[H]
\centering
\begin{tabularx}{0.8\textwidth}{lX}
\toprule
\textbf{Aspek} & \textbf{Karakteristik} \\
\midrule
Usia & 30-55 tahun \\
Usaha & Warung makan, kafe, katering \\
Teknologi & Cukup mahir menggunakan smartphone \\
Kebutuhan & Mengurangi food waste, menambah pendapatan \\
Pain Point & Kesulitan prediksi permintaan, tidak ada platform khusus \\
\bottomrule
\end{tabularx}
\end{table}

\subsubsection{Persona 2: Customer (Draft)}
\begin{table}[H]
\centering
\begin{tabularx}{0.8\textwidth}{lX}
\toprule
\textbf{Aspek} & \textbf{Karakteristik} \\
\midrule
Usia & 18-40 tahun (milenial dan Gen Z) \\
Pekerjaan & Mahasiswa, karyawan swasta \\
Kebutuhan & Makanan berkualitas dengan harga terjangkau \\
Pain Point & Ketidakyakinan terhadap kelayakan makanan diskon \\
\bottomrule
\end{tabularx}
\end{table}

\subsubsection{Persona 3: Admin Platform (Draft)}
\begin{table}[H]
\centering
\begin{tabularx}{0.8\textwidth}{lX}
\toprule
\textbf{Aspek} & \textbf{Karakteristik} \\
\midrule
Usia & 25-45 tahun \\
Peran & Operator platform, moderator \\
Kebutuhan & Monitoring transaksi, verifikasi UMKM \\
\bottomrule
\end{tabularx}
\end{table}

\subsection{Rencana Desain UI/UX}

\textbf{STATUS}: Belum dimulai. Desain akan dibuat setelah user research selesai.

\subsubsection{Tools yang Akan Digunakan}
\begin{itemize}[leftmargin=*]
    \item \textbf{Wireframing} — Figma atau Adobe XD
    \item \textbf{Prototyping} — Figma Interactive Prototype
    \item \textbf{Design System} — Akan dibuat custom design system
\end{itemize}

\subsubsection{Halaman yang Akan Di-design}
\begin{enumerate}[leftmargin=*]
    \item \textbf{Customer Flow}
    \begin{itemize}
        \item Landing page
        \item Marketplace (list produk)
        \item Detail produk
        \item Checkout dan payment
        \item Order tracking
        \item Profile dan impact dashboard
    \end{itemize}
    
    \item \textbf{UMKM Flow}
    \begin{itemize}
        \item Dashboard overview
        \item Tambah produk + Food Eligibility Assessment
        \item Kelola produk dan inventory
        \item Riwayat transaksi
        \item Food Rescue Score dashboard
    \end{itemize}
    
    \item \textbf{Admin Flow}
    \begin{itemize}
        \item Admin dashboard
        \item User management
        \item UMKM verification
        \item Content moderation
    \end{itemize}
\end{enumerate}

\subsection{Prinsip Desain yang Akan Diterapkan}

\begin{enumerate}[leftmargin=*]
    \item \textbf{Konsistensi} — Penggunaan warna, tipografi, dan komponen UI yang konsisten.
    
    \item \textbf{Hierarki Visual} — Informasi penting (harga diskon, sisa waktu) akan ditampilkan dengan ukuran yang menonjol.
    
    \item \textbf{Feedback} — Setiap aksi pengguna akan mendapatkan konfirmasi visual.
    
    \item \textbf{Mobile-First} — Desain akan dimulai dari layar mobile.
\end{enumerate}

\subsection{Rencana Accessibility (Target WCAG 2.2 Level AA)}

\begin{itemize}[leftmargin=*]
    \item \textbf{Color Contrast} — Akan dipatuhi rasio kontras minimal 4.5:1 untuk teks.
    
    \item \textbf{Keyboard Navigation} — Seluruh fitur utama akan dapat diakses menggunakan keyboard.
    
    \item \textbf{Alt Text} — Semua gambar informational akan memiliki deskripsi alternatif.
    
    \item \textbf{Responsive Design} — Akan optimal di berbagai ukuran layar (mobile 360px+, desktop 1024px+).
    
    \item \textbf{Font Size} — Ukuran font minimal 14px untuk mobile, 16px untuk desktop.
    
    \item \textbf{ARIA Attributes} — Akan digunakan untuk meningkatkan aksesibilitas komponen interaktif.
\end{itemize}

\subsection{Timeline Desain UI/UX}

\begin{table}[H]
\centering
\caption{Rencana Timeline Desain}
\label{tab:design_timeline}
\begin{tabularx}{\textwidth}{lXl}
\toprule
\textbf{Fase} & \textbf{Aktivitas} & \textbf{Durasi} \\
\midrule
Research & User interviews dan survey & 1-2 minggu \\
Wireframing & Low-fidelity wireframes & 1 minggu \\
Design & High-fidelity mockups & 2 minggu \\
Prototyping & Interactive prototype & 1 minggu \\
Testing & Usability testing & 1 minggu \\
Refinement & Iterasi berdasarkan feedback & 1 minggu \\
\bottomrule
\end{tabularx}
\end{table}

\textbf{TOTAL ESTIMASI}: 7-9 minggu untuk complete UI/UX design.

\subsection{Deliverables Desain}

Hasil akhir dari fase desain UI/UX akan mencakup:

\begin{itemize}[leftmargin=*]
    \item \textbf{User Personas} — Dokumen lengkap untuk 3 target user
    \item \textbf{User Journey Maps} — Flow lengkap untuk setiap persona
    \item \textbf{Wireframes} — Low-fidelity untuk semua halaman utama
    \item \textbf{High-Fidelity Mockups} — Desain visual final untuk development
    \item \textbf{Interactive Prototype} — Clickable prototype untuk testing
    \item \textbf{Design System} — Component library dan style guide
    \item \textbf{Accessibility Checklist} — Dokumentasi compliance WCAG 2.2
\end{itemize}

\textbf{CATATAN PENTING}: Semua deliverables di atas adalah RENCANA yang akan dikerjakan pada fase development. Saat ini belum ada mockup, wireframe, atau design system yang sudah dibuat.
